\begin{lexlean}{IdsRT}
\useglossary{lexlean.std.nat@1.2.0}
\useglossary{spec.titles@1.0.0}
\importmodule{Rows}
\title{Runtime conformance claims}
\begin{semanticmodule}
definition{
  name: rt01,
  parameters: [],
  result: Rows::IdRow,
  body: (record Rows::IdRow {id := "RT-01", level := "build", suite := "runtime", statement := "The runtime of language 1.2 is LexLean source: every module of its registry is one native core module of the runtime's own project, which holds no file of Lean source, the registry lists only the modules in dependency order, every declaration pins its exact axioms, and a source that is missing, weakened to an inexact policy, or imports anything but `Init` and an earlier runtime module is refused; the one-time migration input is absent, no native source names a runtime copy inside a module, and no Lean file declares a runtime namespace but the released output of a language-1.0 or 1.1 project.", sections := "§17.12"}),
  executable: true
}

definition{
  name: rt02,
  parameters: [],
  result: Rows::IdRow,
  body: (record Rows::IdRow {id := "RT-02", level := "build", suite := "runtime", statement := "The runtime text of languages 1.0 and 1.1 is the one declared exception: the portable runtime and the decoder of core module 1 are the released bytes, held as data under `language/runtime-1.2/frozen/`, read by the language-1.0 and 1.1 backends alone, excluded from the language-1.0 and 1.1 identities, and reproduced byte for byte by the released language-1.1 and 1.0 output; a one-byte change of either is not the released output.", sections := "§17.12, §21.2"}),
  executable: true
}

definition{
  name: rt03,
  parameters: [],
  result: Rows::IdRow,
  body: (record Rows::IdRow {id := "RT-03", level := "build", suite := "runtime", statement := "A verification compiles every runtime module that a project's generated modules, or the calculus modules its certificates are checked against, import before the project's modules, replays it with `leanchecker`, and audits each of its declarations against the axioms its row pins; the committed verification of a project that uses the collections prints them, the runtime's own project verifies under the pinned Lean, and a row that pins another axiom set than its declaration observes fails with LLV7005.", sections := "§17.12, §22.3, §22.5, §22.6"}),
  executable: true
}

definition{
  name: rt04,
  parameters: [],
  result: Rows::IdRow,
  body: (record Rows::IdRow {id := "RT-04", level := "build", suite := "runtime", statement := "No Lean runtime text remains in the language-1.2 compiler: a generated language-1.2 module imports the runtime modules whose constants it names and declares nothing under a runtime namespace, and the only Lean runtime text in the compiler's source outside unit-test items, in one string literal or in several written together, is the decoder of core module 2, which loads native modules and so cannot be one; a runtime text planted in any of its shapes, raw or escaped, with or without a namespace line, split across literals, or after a test module, is found.", sections := "§17.12, §18.1"}),
  executable: true
}
\end{semanticmodule}
\end{lexlean}
